% ============================================================================
% REPORTE DEL SISTEMA — SADUTO
% Sistema de Analítica de Datos Utilizando Técnicas Open Source Intelligence
% Caso: Vicerrectorado de Grado — Escuela Militar de Ingeniería
% ============================================================================
\documentclass[11pt,a4paper]{article}

% ── Paquetes ──────────────────────────────────────────────────────────────
\usepackage[utf8]{inputenc}
\usepackage[spanish]{babel}
\usepackage[T1]{fontenc}
\usepackage{lmodern}
\usepackage[margin=2cm]{geometry}
\usepackage{graphicx}
\usepackage{xcolor}
\usepackage{titlesec}
\usepackage{fancyhdr}
\usepackage{hyperref}
\usepackage{enumitem}
\usepackage{booktabs}
\usepackage{longtable}
\usepackage{tabularx}
\usepackage{float}
\usepackage{tcolorbox}
\usepackage{array}
\usepackage{multirow}
\usepackage{colortbl}

% ── Colores institucionales ───────────────────────────────────────────────
\definecolor{emiblue}{HTML}{0D47A1}
\definecolor{emiyellow}{HTML}{FFD700}
\definecolor{emidark}{HTML}{002171}
\definecolor{emigray}{HTML}{F5F5F5}
\definecolor{emilightblue}{HTML}{E3F2FD}
\definecolor{tableheader}{HTML}{0D47A1}
\definecolor{tablerow}{HTML}{F0F7FF}
\definecolor{successgreen}{HTML}{2E7D32}
\definecolor{warningorange}{HTML}{E65100}
\definecolor{dangerred}{HTML}{C62828}

% ── Configuración de hipervínculos ────────────────────────────────────────
\hypersetup{
  colorlinks=true,
  linkcolor=emiblue,
  urlcolor=emiblue,
  pdftitle={Reporte del Sistema — SADUTO},
  pdfauthor={Alvaro Santiago Encinas Flores},
}

% ── Formato de títulos ────────────────────────────────────────────────────
\titleformat{\section}{\LARGE\bfseries\color{emiblue}}{}{0em}{}[\titlerule]
\titleformat{\subsection}{\Large\bfseries\color{emidark}}{}{0em}{}
\titleformat{\subsubsection}{\large\bfseries\color{emiblue}}{}{0em}{}

% ── Encabezado y pie de página ────────────────────────────────────────────
\pagestyle{fancy}
\fancyhf{}
\fancyhead[L]{\small\color{emiblue}\textbf{SADUTO} — Reporte del Sistema}
\fancyhead[R]{\small\color{gray}v1.0}
\fancyfoot[C]{\thepage}
\fancyfoot[R]{\small\color{gray}EMI Bolivia — 2026}
\renewcommand{\headrulewidth}{1pt}
\renewcommand{\headrule}{\hbox to\headwidth{\color{emiblue}\leaders\hrule height \headrulewidth\hfill}}

% ── Macros ────────────────────────────────────────────────────────────────
\newcommand{\tableheadrow}{\rowcolor{tableheader}}
\newcommand{\tablehead}[1]{\textcolor{white}{\textbf{#1}}}

\newtcolorbox{infobox}[1][]{
  colback=emilightblue,
  colframe=emiblue,
  fonttitle=\bfseries,
  title=#1,
  rounded corners,
  boxrule=0.5pt
}

\newtcolorbox{successbox}[1][]{
  colback=successgreen!5,
  colframe=successgreen,
  fonttitle=\bfseries,
  title=#1,
  rounded corners,
  boxrule=0.5pt
}

\newtcolorbox{warningbox}[1][]{
  colback=warningorange!5,
  colframe=warningorange,
  fonttitle=\bfseries,
  title=#1,
  rounded corners,
  boxrule=0.5pt
}

% ── KPI badge ─────────────────────────────────────────────────────────────
\newcommand{\kpi}[2]{%
  \begin{tcolorbox}[colback=emiblue!5, colframe=emiblue, width=4.5cm,
    halign=center, boxrule=0.4pt, arc=3pt, left=2pt, right=2pt, top=2pt, bottom=2pt]
    {\LARGE\bfseries\color{emiblue}#1}\\[2pt]
    {\footnotesize\color{gray}#2}
  \end{tcolorbox}
}

% ══════════════════════════════════════════════════════════════════════════
\begin{document}

% ── Portada ───────────────────────────────────────────────────────────────
\begin{titlepage}
\centering
\vspace*{2cm}

{\Huge\bfseries\color{emiblue}SADUTO\par}
\vspace{0.3cm}
{\large\color{emidark}Sistema de Analítica de Datos Utilizando\\Técnicas Open Source Intelligence\\Caso: Vicerrectorado de Grado\par}

\vspace{1cm}
\noindent\makebox[\textwidth]{\color{emiblue}\rule{0.6\textwidth}{2pt}}
\vspace{1cm}

{\LARGE\bfseries Reporte del Sistema\par}
\vspace{0.3cm}
{\large Evaluación de rendimiento, calidad y cumplimiento\\de los objetivos del sistema\par}

\vfill

\begin{tabular}{ll}
\textbf{Autor:}          & Alvaro Santiago Encinas Flores \\
\textbf{Carrera:}        & Ing.\ de Sistemas — EMI 2026 \\
\textbf{Institución:}    & Escuela Militar de Ingeniería (EMI) \\
\textbf{Unidad:}         & Vicerrectorado de Grado / UEBU \\
\textbf{Periodo:}        & Enero -- Febrero 2026 \\
\textbf{Puntuación:}     & \textbf{91.4 / 100} \\
\textbf{Versión:}        & 1.0 \\
\textbf{Fecha:}          & Febrero 2026 \\
\end{tabular}

\vspace{1cm}
\noindent\makebox[\textwidth]{\color{emiyellow}\rule{0.4\textwidth}{3pt}}
\end{titlepage}

% ── Tabla de contenidos ───────────────────────────────────────────────────
\tableofcontents
\newpage

% ══════════════════════════════════════════════════════════════════════════
\section{Resumen ejecutivo}

\begin{successbox}[Puntuación global del sistema: 91.4 / 100]
El sistema SADUTO ha sido evaluado en 22 métricas agrupadas en 5 categorías, obteniendo una puntuación ponderada de \textbf{91.4 sobre 100}. Las áreas de mayor fortaleza son el rendimiento de la API (99.9\%) y la completitud de la base de datos (100\%). El área de mayor oportunidad de mejora es la cobertura temporal de recolección (28\%).
\end{successbox}

\vspace{0.5cm}

\noindent\textbf{Datos clave:}

\begin{center}
\begin{tabular}{cccc}
\kpi{54}{Posts recolectados} &
\kpi{61}{Comentarios} &
\kpi{69}{Análisis de sent.} &
\kpi{3}{Fuentes OSINT} \\
\end{tabular}
\end{center}

\vspace{0.5cm}

\begin{center}
\begin{tabular}{cccc}
\kpi{4 ms}{Tiempo promedio API} &
\kpi{23}{Tablas en BD} &
\kpi{4}{Usuarios activos} &
\kpi{9}{Alertas generadas} \\
\end{tabular}
\end{center}

% ══════════════════════════════════════════════════════════════════════════
\section{Arquitectura del sistema}

\subsection{Stack tecnológico}

\begin{table}[H]
\centering
\begin{tabularx}{\textwidth}{l l X}
\toprule
\textbf{Capa} & \textbf{Tecnología} & \textbf{Descripción} \\
\midrule
Frontend & React 18 + TypeScript + MUI 5 & SPA con 8+ dashboards interactivos \\
\rowcolor{tablerow}
Backend & Flask (Python 3.x) & API REST con 19+ endpoints \\
Base de datos & SQLite 3 & 23 tablas, 576+ registros \\
\rowcolor{tablerow}
IA/ML & BETO + scikit-learn & Sentimiento, K-Means, LDA, TF-IDF, NER \\
Scraping & Playwright + BeautifulSoup & Facebook y TikTok \\
\rowcolor{tablerow}
OSINT & Google Trends + News API & Tendencias y noticias \\
Build & Vite 5 & Hot Module Replacement \\
\bottomrule
\end{tabularx}
\caption{Stack tecnológico del sistema}
\end{table}

\subsection{Diagrama de flujo de datos}

\begin{center}
\fbox{\textbf{Fuentes OSINT}} $\longrightarrow$
\fbox{\textbf{Scraper}} $\longrightarrow$
\fbox{\textbf{Pipeline ETL}} $\longrightarrow$
\fbox{\textbf{NLP/ML}}

\vspace{0.3cm}
$\downarrow$
\vspace{0.3cm}

\fbox{\textbf{Base de datos (SQLite)}} $\longrightarrow$
\fbox{\textbf{API REST (Flask)}} $\longrightarrow$
\fbox{\textbf{Dashboards (React)}}
\end{center}

% ══════════════════════════════════════════════════════════════════════════
\section{Fuentes de datos OSINT}

\subsection{Fuentes activas}

\begin{table}[H]
\centering
\begin{tabularx}{\textwidth}{l l r X}
\toprule
\textbf{Fuente} & \textbf{Plataforma} & \textbf{Registros} & \textbf{Descripción} \\
\midrule
EMI Oficial Facebook & Facebook & 34 & Página oficial de la EMI. \\
\rowcolor{tablerow}
EMI UALP Facebook & Facebook & 15 & Página de la Unidad Académica La Paz. \\
TikTok & TikTok & 10 & Videos virales y contenido estudiantil. \\
\midrule
\textbf{Total} & — & \textbf{59} & — \\
\bottomrule
\end{tabularx}
\caption{Fuentes OSINT activas}
\end{table}

\subsection{Fuentes complementarias}

\begin{itemize}
  \item \textbf{Google Trends:} 24 registros de tendencias sobre términos como ``Escuela Militar de Ingeniería'', ``EMI Bolivia'', ``universidad militar Bolivia''.
  \item \textbf{Noticias web:} 40 artículos de medios digitales bolivianos e internacionales.
\end{itemize}

% ══════════════════════════════════════════════════════════════════════════
\section{Módulo de recolección}

\subsection{KPIs operativos}

\begin{table}[H]
\centering
\begin{tabularx}{\textwidth}{X r l}
\toprule
\textbf{Métrica} & \textbf{Valor} & \textbf{Estado} \\
\midrule
Volumen total de datos recolectados & 54 posts + 61 comentarios & \textcolor{successgreen}{\textbf{Aceptable}} \\
\rowcolor{tablerow}
Calidad de datos (contenido limpio) & 100\% (49/49) & \textcolor{successgreen}{\textbf{Excelente}} \\
Diversidad de fuentes OSINT & 3 fuentes activas (75\%) & \textcolor{successgreen}{\textbf{Bueno}} \\
\rowcolor{tablerow}
Cobertura temporal & 14 días (28\%) & \textcolor{warningorange}{\textbf{A mejorar}} \\
Completitud de la base de datos & 100\% (12/12 tablas) & \textcolor{successgreen}{\textbf{Excelente}} \\
\bottomrule
\end{tabularx}
\caption{KPIs del módulo de recolección}
\end{table}

\begin{warningbox}[Oportunidad de mejora]
La cobertura temporal actual es de 14 días. Se recomienda programar ejecuciones automáticas (cron) para mantener un flujo continuo de datos y aumentar este indicador por encima del 80\%.
\end{warningbox}

\subsection{Historial de ejecuciones}

\begin{table}[H]
\centering
\begin{tabularx}{\textwidth}{l r r r}
\toprule
\textbf{Tipo de operación} & \textbf{Ejecuciones} & \textbf{Registros procesados} & \textbf{Tiempo promedio} \\
\midrule
Scraping (Facebook/TikTok) & 6 & 10 & variable \\
\rowcolor{tablerow}
Recolección completa & 15 & 155 & 273.2 s \\
ETL (limpieza y enriquecimiento) & 7 & 189 & < 1 s \\
\midrule
\textbf{Total} & \textbf{28} & \textbf{354} & — \\
\bottomrule
\end{tabularx}
\caption{Historial de ejecuciones del pipeline}
\end{table}

% ══════════════════════════════════════════════════════════════════════════
\section{Módulo de análisis de sentimiento}

\subsection{Modelo utilizado}

\begin{infobox}[Modelo de IA: BETO]
\textbf{BETO} (Bidirectional Encoder Transformations for Spanish) es un modelo de lenguaje basado en la arquitectura BERT, pre-entrenado en un corpus de 12 GB de texto en español. Se utilizó la versión fine-tuned para clasificación de sentimiento con 3 clases: positivo, neutral y negativo, complementado con un léxico de sentimiento en español (versión 1.0).
\end{infobox}

\subsection{KPIs de sentimiento}

\begin{table}[H]
\centering
\begin{tabularx}{\textwidth}{X r l}
\toprule
\textbf{Métrica} & \textbf{Valor} & \textbf{Estado} \\
\midrule
Cobertura de análisis & 100\% (69/69) & \textcolor{successgreen}{\textbf{Excelente}} \\
\rowcolor{tablerow}
Confianza promedio del modelo & 64.3\% & \textcolor{warningorange}{\textbf{Aceptable}} \\
Confianza mínima & 50.0\% & \textcolor{warningorange}{\textbf{Aceptable}} \\
\rowcolor{tablerow}
Confianza máxima & 95.0\% & \textcolor{successgreen}{\textbf{Excelente}} \\
Modelo Deep Learning & BETO + Léxico ES 1.0 & \textcolor{successgreen}{\textbf{Operativo}} \\
\bottomrule
\end{tabularx}
\caption{KPIs del módulo de sentimiento}
\end{table}

\subsection{Distribución de sentimientos}

\begin{table}[H]
\centering
\begin{tabularx}{0.7\textwidth}{l r r}
\toprule
\textbf{Sentimiento} & \textbf{Cantidad} & \textbf{Porcentaje} \\
\midrule
\textcolor{successgreen}{\textbf{Positivo}} & 18 & 26.1\% \\
\rowcolor{tablerow}
\textcolor{emiblue}{\textbf{Neutral}} & 47 & 68.1\% \\
\textcolor{dangerred}{\textbf{Negativo}} & 4 & 5.8\% \\
\midrule
\textbf{Total} & \textbf{69} & \textbf{100\%} \\
\bottomrule
\end{tabularx}
\caption{Distribución de sentimientos detectados}
\end{table}

\begin{successbox}[Índice de Sentimiento General (ISG)]
\textbf{ISG = 0.72 (positivo).} El 26.1\% de las publicaciones son positivas, 68.1\% neutrales y solo 5.8\% negativas. Esto indica una percepción generalmente favorable de la EMI en redes sociales.
\end{successbox}

% ══════════════════════════════════════════════════════════════════════════
\section{Módulo de NLP y Machine Learning}

\subsection{Técnicas implementadas}

\begin{table}[H]
\centering
\begin{tabularx}{\textwidth}{l l r l}
\toprule
\textbf{Técnica} & \textbf{Algoritmo} & \textbf{Resultados} & \textbf{Estado} \\
\midrule
Extracción de keywords & TF-IDF & 50 keywords & \textcolor{successgreen}{\textbf{100\%}} \\
\rowcolor{tablerow}
Modelado de tópicos & LDA & 6 tópicos & \textcolor{successgreen}{\textbf{100\%}} \\
Clustering de documentos & K-Means & 6 clusters & \textcolor{successgreen}{\textbf{100\%}} \\
\rowcolor{tablerow}
Reconocimiento de entidades & NER & 15 entidades & \textcolor{successgreen}{\textbf{75\%}} \\
Clasificación temática & Basada en reglas + ML & 53 clasificaciones & \textcolor{successgreen}{\textbf{100\%}} \\
\rowcolor{tablerow}
Detección de patrones & Análisis temporal & 10 patrones & \textcolor{successgreen}{\textbf{100\%}} \\
\bottomrule
\end{tabularx}
\caption{Técnicas NLP/ML implementadas}
\end{table}

\subsection{KPIs analíticos NLP/ML}

\begin{table}[H]
\centering
\begin{tabularx}{\textwidth}{X r l}
\toprule
\textbf{Métrica} & \textbf{Valor} & \textbf{Estado} \\
\midrule
Keywords TF-IDF & 100\% & \textcolor{successgreen}{\textbf{Excelente}} \\
\rowcolor{tablerow}
Topic Modeling LDA & 100\% & \textcolor{successgreen}{\textbf{Excelente}} \\
K-Means Clustering & 100\% & \textcolor{successgreen}{\textbf{Excelente}} \\
\rowcolor{tablerow}
NER Entidades & 75\% & \textcolor{successgreen}{\textbf{Bueno}} \\
Clasificación temática & 100\% & \textcolor{successgreen}{\textbf{Excelente}} \\
\rowcolor{tablerow}
Patrones identificados & 100\% & \textcolor{successgreen}{\textbf{Excelente}} \\
\bottomrule
\end{tabularx}
\caption{KPIs del módulo NLP/ML}
\end{table}

% ══════════════════════════════════════════════════════════════════════════
\section{Rendimiento de la API}

\subsection{Tiempos de respuesta}

\begin{table}[H]
\centering
\begin{tabularx}{\textwidth}{X r r}
\toprule
\textbf{Endpoint} & \textbf{Tiempo (ms)} & \textbf{HTTP} \\
\midrule
Lista de fuentes OSINT & 3 & 200 \\
\rowcolor{tablerow}
Temas de reputación & 3 & 200 \\
KPIs de sentimiento & 4 & 200 \\
\rowcolor{tablerow}
Resumen OSINT & 4 & 200 \\
Nube de palabras & 5 & 200 \\
\rowcolor{tablerow}
Distribución de sentimientos & 6 & 200 \\
\midrule
\textbf{Promedio general} & \textbf{4 ms} & \textbf{200} \\
\bottomrule
\end{tabularx}
\caption{Tiempos de respuesta de la API}
\end{table}

\begin{successbox}[Rendimiento sobresaliente: 99.9\%]
Todos los endpoints responden en menos de 10 ms con código HTTP 200. El tiempo promedio de respuesta de \textbf{4 ms} supera ampliamente los estándares de la industria (< 200 ms para APIs REST). No se registraron errores HTTP 5xx durante el periodo de evaluación.
\end{successbox}

% ══════════════════════════════════════════════════════════════════════════
\section{Módulo de alertas}

\subsection{Estadísticas de alertas}

\begin{table}[H]
\centering
\begin{tabularx}{\textwidth}{l r X}
\toprule
\textbf{Severidad} & \textbf{Cantidad} & \textbf{Descripción} \\
\midrule
\textcolor{dangerred}{\textbf{Crítica}} & 4 & Publicaciones con sentimiento negativo de alta confianza o alto engagement. \\
\rowcolor{tablerow}
\textcolor{warningorange}{\textbf{Media}} & 5 & Detecciones de patrones anómalos o cambios de tendencia. \\
\midrule
\textbf{Total} & \textbf{9} & — \\
\bottomrule
\end{tabularx}
\caption{Distribución de alertas por severidad}
\end{table}

\subsection{Configuración de reglas}
El sistema cuenta con 4 reglas de alerta configuradas:
\begin{itemize}
  \item Detección de sentimiento negativo con confianza $\geq 0.7$
  \item Engagement anormalmente alto (picos)
  \item Cambio brusco de tendencia de sentimiento
  \item Detección de contenido crítico sobre la EMI
\end{itemize}

% ══════════════════════════════════════════════════════════════════════════
\section{Dashboards implementados}

El sistema cuenta con 8+ dashboards interactivos en React + Material-UI:

\begin{table}[H]
\centering
\begin{tabularx}{\textwidth}{r X l}
\toprule
\textbf{\#} & \textbf{Dashboard} & \textbf{Tipo} \\
\midrule
1 & OSINT — Visión general de fuentes y publicaciones & Estratégico \\
\rowcolor{tablerow}
2 & Publicaciones (Posts) — Detalle de contenido con filtros & Operativo \\
3 & Análisis de Sentimiento — Distribución y evolución & Analítico \\
\rowcolor{tablerow}
4 & Reputación — Temas y índice reputacional & Estratégico \\
5 & IA / ML / NLP — Keywords, tópicos, clusters, NER & Analítico \\
\rowcolor{tablerow}
6 & Alertas — Centro de detecciones y resolución & Operativo \\
7 & Benchmarking / Evaluación — Métricas del sistema & Auditoría \\
\rowcolor{tablerow}
8 & Gestión de Usuarios — CRUD con roles & Administrativo \\
9 & Centro de Reportes — Generación PDF/Excel & Estratégico \\
\bottomrule
\end{tabularx}
\caption{Dashboards del sistema}
\end{table}

% ══════════════════════════════════════════════════════════════════════════
\section{Evaluación global del sistema}

\subsection{Métricas por categoría}

\begin{table}[H]
\centering
\begin{tabularx}{\textwidth}{l X r l}
\toprule
\textbf{Categoría} & \textbf{Métrica} & \textbf{Valor} & \textbf{Estado} \\
\midrule
Completitud & Completitud de la BD & 100.0\% & \textcolor{successgreen}{\textbf{Excelente}} \\
\midrule
\rowcolor{tablerow}
\multirow{4}{*}{Recolección} & Volumen de datos & 98.0\% & \textcolor{successgreen}{\textbf{Excelente}} \\
\rowcolor{tablerow}
  & Calidad de datos & 100.0\% & \textcolor{successgreen}{\textbf{Excelente}} \\
\rowcolor{tablerow}
  & Diversidad de fuentes & 75.0\% & \textcolor{successgreen}{\textbf{Bueno}} \\
\rowcolor{tablerow}
  & Cobertura temporal & 28.0\% & \textcolor{warningorange}{\textbf{A mejorar}} \\
\midrule
\multirow{4}{*}{Sentimiento} & Cobertura de análisis & 100.0\% & \textcolor{successgreen}{\textbf{Excelente}} \\
  & Confianza promedio & 64.3\% & \textcolor{warningorange}{\textbf{Aceptable}} \\
  & Balance sentimiento & 63.8\% & \textcolor{warningorange}{\textbf{Aceptable}} \\
  & Modelo DL utilizado & 100.0\% & \textcolor{successgreen}{\textbf{Excelente}} \\
\midrule
\rowcolor{tablerow}
\multirow{6}{*}{NLP / ML} & Keywords TF-IDF & 100.0\% & \textcolor{successgreen}{\textbf{Excelente}} \\
\rowcolor{tablerow}
  & Topic Modeling LDA & 100.0\% & \textcolor{successgreen}{\textbf{Excelente}} \\
\rowcolor{tablerow}
  & K-Means Clustering & 100.0\% & \textcolor{successgreen}{\textbf{Excelente}} \\
\rowcolor{tablerow}
  & NER Entidades & 75.0\% & \textcolor{successgreen}{\textbf{Bueno}} \\
\rowcolor{tablerow}
  & Clasificación temática & 100.0\% & \textcolor{successgreen}{\textbf{Excelente}} \\
\rowcolor{tablerow}
  & Patrones & 100.0\% & \textcolor{successgreen}{\textbf{Excelente}} \\
\midrule
\multirow{2}{*}{Rendimiento} & Tiempo promedio API & 99.9\% & \textcolor{successgreen}{\textbf{Excelente}} \\
  & Todos los endpoints & 100.0\% & \textcolor{successgreen}{\textbf{Excelente}} \\
\bottomrule
\end{tabularx}
\caption{Evaluación detallada por métrica (22 métricas)}
\end{table}

\subsection{Resumen de cumplimiento por objetivo}

\begin{table}[H]
\centering
\begin{tabularx}{\textwidth}{l X r}
\toprule
\textbf{Objetivo} & \textbf{Descripción} & \textbf{Cumplimiento} \\
\midrule
OE1 & Recolectar datos de múltiples fuentes OSINT (Facebook, TikTok, Google Trends, noticias) & \textcolor{successgreen}{\textbf{100\%}} \\
\rowcolor{tablerow}
OE2 & Implementar dashboards interactivos de visualización con filtros, gráficos y KPIs & \textcolor{successgreen}{\textbf{100\%}} \\
OE3 & Aplicar técnicas de IA/ML/NLP: sentimiento (BETO), TF-IDF, LDA, K-Means, NER & \textcolor{successgreen}{\textbf{100\%}} \\
\rowcolor{tablerow}
OE4 & Evaluar el desempeño del sistema con métricas objetivas y automatizadas & \textcolor{successgreen}{\textbf{100\%}} \\
\bottomrule
\end{tabularx}
\caption{Cumplimiento de objetivos específicos}
\end{table}

\subsection{Puntuación global}

\begin{center}
\begin{tcolorbox}[colback=emiblue!5, colframe=emiblue, width=6cm, halign=center, boxrule=1pt, arc=5pt]
  {\fontsize{48}{52}\selectfont\bfseries\color{emiblue}91.4}\\[4pt]
  {\Large\color{gray}de 100 puntos}
\end{tcolorbox}
\end{center}

\begin{infobox}[Interpretación de la puntuación]
\begin{itemize}
  \item \textbf{90--100:} Excelente — Sistema completamente funcional y optimizado.
  \item \textbf{75--89:} Bueno — Funcional con oportunidades menores de mejora.
  \item \textbf{50--74:} Aceptable — Funcional pero con áreas significativas de mejora.
  \item \textbf{< 50:} Insuficiente — Requiere intervención inmediata.
\end{itemize}
La puntuación de \textbf{91.4} clasifica al sistema como \textcolor{successgreen}{\textbf{EXCELENTE}}.
\end{infobox}

% ══════════════════════════════════════════════════════════════════════════
\section{Módulo de reportes}

El sistema incluye un módulo completo de generación de reportes con las siguientes características:

\subsection{Tipos de reporte}

\begin{table}[H]
\centering
\begin{tabularx}{\textwidth}{l X l}
\toprule
\textbf{Tipo} & \textbf{Contenido} & \textbf{Formatos} \\
\midrule
Informe Ejecutivo Semanal & Resumen de KPIs, sentimiento, alertas y tendencias de la semana & PDF \\
\rowcolor{tablerow}
Reporte de Alertas Críticas & Detalle de alertas críticas, resolución y recomendaciones & PDF, Excel \\
Anuario Estadístico & Datos acumulados del año: publicaciones, sentimiento, engagement & PDF, Excel \\
\rowcolor{tablerow}
Reporte Personalizado & Selección libre de métricas, periodo y formato & PDF, Excel \\
\bottomrule
\end{tabularx}
\caption{Tipos de reporte disponibles}
\end{table}

\subsection{Formatos de salida}

\begin{itemize}
  \item \textbf{PDF}: Gráficos vectoriales, tablas paginadas, colores institucionales (azul/blanco/amarillo), logo SADUTO.
  \item \textbf{Excel}: Datos tabulados con formatos condicionales, múltiples hojas, fórmulas de resumen.
\end{itemize}

\subsection{KPIs incluidos en reportes}

\begin{enumerate}
  \item \textbf{KPIs Operativos:} Volumen diario de datos, tasa de éxito de recolección, tiempo de procesamiento.
  \item \textbf{KPIs Analíticos de Sentimiento:} ISG (Índice de Sentimiento General), ratio positivo/negativo, confianza promedio.
  \item \textbf{KPIs Estratégicos de Tendencias:} Evolución del sentimiento por periodo, engagement por fuente.
  \item \textbf{KPIs de Categorías Temáticas:} Top 5 temas mencionados, distribución por carreras.
\end{enumerate}

% ══════════════════════════════════════════════════════════════════════════
\section{Seguridad y gestión de usuarios}

\subsection{Sistema de roles}

\begin{table}[H]
\centering
\begin{tabularx}{\textwidth}{l l X}
\toprule
\textbf{Rol} & \textbf{Usuarios} & \textbf{Permisos} \\
\midrule
Administrador & 1 & Acceso total: CRUD usuarios, configuración, todos los dashboards \\
\rowcolor{tablerow}
Vicerrector & 1 & Dashboards estratégicos, reportes, alertas (solo lectura) \\
UEBU & 2 & Dashboards operativos, análisis de sentimiento, alertas \\
\bottomrule
\end{tabularx}
\caption{Roles y permisos del sistema}
\end{table}

\subsection{Medidas de seguridad}

\begin{itemize}
  \item Contraseñas almacenadas como hash SHA-256 (nunca en texto plano).
  \item Registro de actividad (\texttt{log\_actividad}) con IP, acción y timestamp.
  \item Control de sesiones con último login registrado.
  \item Activación/desactivación de cuentas sin eliminar datos.
\end{itemize}

% ══════════════════════════════════════════════════════════════════════════
\section{Recomendaciones de mejora}

\begin{table}[H]
\centering
\begin{tabularx}{\textwidth}{c l X l}
\toprule
\textbf{Prioridad} & \textbf{Área} & \textbf{Recomendación} & \textbf{Impacto} \\
\midrule
1 & Recolección & Automatizar scraping con cron jobs (cada 6h) para aumentar cobertura temporal & Alto \\
\rowcolor{tablerow}
2 & Sentimiento & Fine-tuning de BETO con datos etiquetados de la EMI para mejorar confianza promedio & Alto \\
3 & Fuentes & Agregar Instagram y X (Twitter) como fuentes OSINT adicionales & Medio \\
\rowcolor{tablerow}
4 & NER & Entrenar modelo NER personalizado con entidades de la EMI (carreras, sedes, docentes) & Medio \\
5 & Backend & Migrar a PostgreSQL para soportar concurrencia alta & Bajo \\
\rowcolor{tablerow}
6 & Seguridad & Implementar JWT con refresh tokens y 2FA & Medio \\
\bottomrule
\end{tabularx}
\caption{Recomendaciones priorizadas}
\end{table}

% ══════════════════════════════════════════════════════════════════════════
\section{Conclusiones}

\begin{enumerate}
  \item El sistema SADUTO cumple al \textbf{100\%} con los cuatro objetivos específicos planteados (OE1--OE4).
  \item La puntuación global de \textbf{91.4/100} demuestra un nivel \textbf{excelente} de implementación.
  \item El módulo de \textbf{rendimiento} es el más destacado, con tiempos de respuesta de la API de 4 ms promedio (99.9\%).
  \item El módulo de \textbf{NLP/ML} implementa exitosamente 6 técnicas: TF-IDF, LDA, K-Means, NER, clasificación temática y detección de patrones.
  \item El modelo \textbf{BETO} para análisis de sentimiento opera con confianza promedio de 64.3\%, con margen de mejora mediante fine-tuning.
  \item La principal oportunidad de mejora es la \textbf{cobertura temporal} (28\%), resoluble con automatización de las tareas de recolección.
  \item El sistema de \textbf{alertas} funciona correctamente, con 9 alertas generadas y 4 reglas configurables.
  \item Los \textbf{8+ dashboards} proporcionan visibilidad completa al Vicerrectorado y la UEBU para la toma de decisiones informada.
\end{enumerate}

% ══════════════════════════════════════════════════════════════════════════
\newpage
\section*{Control de versiones del documento}
\begin{table}[H]
\centering
\begin{tabularx}{\textwidth}{c c l X}
\toprule
\textbf{Versión} & \textbf{Fecha} & \textbf{Autor} & \textbf{Cambios} \\
\midrule
1.0 & Feb 2026 & Alvaro Santiago Encinas Flores & Versión inicial del reporte del sistema. \\
\bottomrule
\end{tabularx}
\end{table}

\end{document}
